\section{Introduction}

Empirical mode decomposition (EMD) is a data-driven way to split a nonstationary signal into a small number of oscillatory components, the intrinsic mode functions (IMFs), with no fixed basis \cite{huang1998}. Electroencephalography (EEG) is nonstationary, so it is a natural test: if simple features of the IMFs separate the classes of an EEG task, that supports EMD as a feature extractor.

The coursework pipeline asked whether simple IMF features separate EEG classes and whether the distance used to compare them matters. Its results were at chance on motor imagery and, on mental arithmetic, no more accurate than always predicting rest, and its clustering scores were high. This paper asks a sharper question: does EMD add anything over a plain spectral baseline, on subjects the model has not seen? It makes three contributions: the corrected baseline, with the bugs and their effects (Sections~\ref{sec:data} and~\ref{sec:baseline}); a comparison of per-channel EMD features with band power under subject-held-out evaluation (Sections~\ref{sec:design} to~\ref{sec:emd}); and an account of what the clustering and silhouette scores did and did not measure (Section~\ref{sec:clustering}).

\emph{Revision note.} The primary mistakes in the original coursework report (spring 2026) were describing its data sets with incorrect counts and drawing conclusions that the data did not support. This revised version (fall 2026) uses the proper counts (from the actual run), corrects three bugs (the EMD residue counted as an IMF and zero-padded, the padded tail windows of EEGMAT, and three EEG MMI subjects recorded at 128 Hz and read as 160 Hz), rewrites the clustering discussion, and adds the study in Sections~\ref{sec:design} to~\ref{sec:emd}.
